%% kapa: exam=1 round=20 year=2009 period=1 form=A subject=civil_law subject_ko=민법 seq=16 number=16
%% source: Q-Net 감정평가사 기출문제 (cst00309) artlSeq=5024178 file=extracted/1008/5024178/2009년 제20회 감정평가사 1차시험 기출문제/1교시.hwp
%% conversion: kordoc 4.24.51 -> tools/gen1.py (auto); review=proofread-sonnet; flags=none
\begin{question}{16}
점유에 관한 설명으로 옳지 \uline{않은} 것은? (다툼이 있으면 판례에 의함)
\begin{choices}
  \item 점유보조자는 점유방해자에 대하여 점유보호청구권을 행사할 수 없다.
  \item 공유자 1인이 공유토지 전부를 점유하는 경우, 특별한 사정이 없는 한 다른 공유자의 지분비율의 범위 내에서는 타주점유이다.
  \item 등기된 부동산에 관하여는 특별한 사정이 없는 한 점유자의 권리에 관한 적법추정이 미치지 않는다.
  \item 점유자는 소유의 의사로 선의ㆍ무과실로 점유한 것으로 추정한다.
  \item 점유권에 기인한 소와 본권에 기인한 소는 서로 영향을 미치지 않는다.
\end{choices}
\end{question}
\answer{4}
